\paragraph{Django}
Django \cite{wikipedia_django} es un framework web escrito en Python. Sus principales características son su fácil curva de aprendizaje, el sistema 
de plantillas para las vistas, un panel de administración incorporado y un ORM (Object-Relational Mapping) integrado. 
Esta última característica permite interactuar con bases de datos relacionales utilizando objetos del lenguaje sin necesidad 
de escribir sentencias SQL. Además, incluye por defecto protección contra ataques CSRF, lo que mejora la seguridad de las 
aplicaciones desarrolladas.

\paragraph{.NET}
.NET \cite{wikipedia_dotnet} es un framework desarrollado por Microsoft. Su característica más destacada es el CLR (Common Language Runtime), 
que permite integrar proyectos escritos en diferentes lenguajes compatibles como C\#, F\# y Visual Basic. 
También incluye un ORM (Entity Framework) para la gestión de bases de datos. A diferencia de Django, no ofrece un panel 
de administración preconfigurado, por lo que es necesario desarrollarlo manualmente. Además, aunque su configuración
inicial es más avanzada, requiere mayor conocimiento técnico.

\paragraph{Spring boot}
Spring Boot \cite{ibm_springboot} es un framework basado en Java que simplifica la creación de aplicaciones web al reducir considerablemente 
la configuración inicial. Esto permite a los desarrolladores crear aplicaciones listas para producción con cambios mínimos. 
Si bien no incluye un panel de administración por defecto, su ecosistema es muy amplio y permite una alta personalización. 
También es compatible con herramientas como Hibernate o JPA para la gestión de bases de datos mediante ORM.


Una vez comparados los frameworks backend disponibles, se ha optado por Django debido a su sencilla curva de aprendizaje y a la inclusión 
de un panel de administración por defecto. Esta funcionalidad resulta especialmente útil en este proyecto, ya que permite 
gestionar fácilmente los videojuegos almacenados en la base de datos desde una interfaz gráfica sin necesidad de desarrollar una.